% F2 -- One real mast, seen with 2G alone and then with 3G added on top.
% Adapted from a schematic in a colleague's report on the same corpus.
\begin{figure}[!htbp]
\centering
\fitwidth{%
\begin{tikzpicture}[font=\small,
  ttl/.style={anchor=south,align=center,font=\footnotesize\bfseries,text=clmodel},
  scap/.style={anchor=north,align=center,font=\scriptsize,text=clrule},
  nm/.style={font=\tiny\ttfamily},
]
% a coverage lobe: {angle}{half-width}{reach}{style}
\def\lobe#1#2#3#4{\draw[#4,rotate=#1] (0,#3) ellipse (#2 and #3);}
% the mast itself: a crossed circle
\def\mast{\draw[clmodel,fill=white,line width=0.7pt] (0,0) circle (0.17);
          \draw[clmodel] (-0.12,-0.12) -- (0.12,0.12);
          \draw[clmodel] (-0.12,0.12) -- (0.12,-0.12);}

% ================= panel (a): 2G only ====================================
\begin{scope}[shift={(3.5,0)}]
  \foreach \aa in {45,135,225,315}{
    \lobe{\aa}{0.62}{1.35}{draw=clbase,fill=clbase!16,line width=0.7pt}}
  \mast
  \node[nm,text=clbase] at (45:1.05)  {BKVCHE1};
  \node[nm,text=clbase] at (135:1.05) {BKVCHE2};
  \node[nm,text=clbase] at (225:1.05) {BKVCHE3};
  \node[nm,text=clbase] at (315:1.05) {BKVCHE4};
  \node[ttl] at (0,2.95) {\panelnum{fig:base-station}{a} One mast carries four 2G cells};
  \node[scap,text width=5.0cm] at (0,-2.95)
    {one installation, four names in the data};
\end{scope}

\draw[clrule!50] (7.0,3.05) -- (7.0,-3.20);

% ================= panel (b): 2G and 3G together =========================
\begin{scope}[shift={(10.5,0)}]
  \foreach \aa in {0,90,180,270}{
    \lobe{\aa}{0.62}{1.35}{draw=clink,fill=clink!13,line width=0.7pt}}
  \foreach \aa in {45,135,225,315}{
    \lobe{\aa}{0.40}{0.82}{draw=clbase,fill=clbase!20,line width=0.7pt}}
  \mast
  \node[nm,text=clink] at (0:1.78)   {UKVCHE101};
  \node[nm,text=clink] at (90:1.78)  {UKVCHE102};
  \node[nm,text=clink] at (180:1.78) {UKVCHE103};
  \node[nm,text=clink] at (270:1.78) {UKVCHE104};
  % the handover that looks like a movement
  \draw[-{Stealth[length=4pt]},clmove,thick]
        (0.38,0.38) to[bend left=25] (1.15,0.22);
  \draw[clmove] (1.24,1.44) -- (0.82,0.82);
  \node[anchor=west,font=\scriptsize,text=clmove,align=left,text width=1.95cm]
        at (1.24,1.72) {a new name,\\ nobody moved};
  \node[ttl] at (0,2.95) {\panelnum{fig:base-station}{b} The same mast also carries 3G: eight cells};
  \node[scap,text width=5.4cm] at (0,-2.95)
    {the smaller lobes are the 2G cells of \panelnum{fig:base-station}{a}};
\end{scope}

% ================= the reading ===========================================
\node[draw=clrule,fill=clsoft,rounded corners=3pt,align=left,inner sep=7pt,
      text width=13.6cm,font=\footnotesize,anchor=north] at (7.0,-3.55)
  {\textbf{Eight names, one place.} A phone lying still at the foot of the mast
   sits inside several lobes at once, and the network can hand it from one to the
   next for reasons of signal alone. Every such handover is written into the data
   as a change of antenna, and looks exactly like a movement.};
\end{tikzpicture}}
\caption{One mast, and all the names it answers to}
\label{fig:base-station}

\end{figure}
